\documentclass[12pt]{report}

\usepackage[utf8x]{inputenc}
\usepackage[T1]{fontenc}
\usepackage[frenchb]{babel}
\usepackage{lmodern}

\usepackage{amsmath,amsthm,amssymb}
\usepackage{graphicx}
\usepackage[colorinlistoftodos]{todonotes}
\usepackage{listings}
\usepackage{color}
\usepackage{verbatim}


\definecolor{mygreen}{rgb}{0,0.6,0}
\definecolor{mygray}{rgb}{0.5,0.5,0.5}
\definecolor{mymauve}{rgb}{0.58,0,0.82}

\lstset{ %
	backgroundcolor=\color{white},   % choose the background color
	basicstyle=\footnotesize,        % size of fonts used for the code
	breaklines=true,                 % automatic line breaking only at whitespace
	captionpos=b,                    % sets the caption-position to bottom
	commentstyle=\color{mygreen},    % comment style
	escapeinside={\%*}{*)},          % if you want to add LaTeX within your code
	keywordstyle=\color{blue},       % keyword style
	stringstyle=\color{mymauve},     % string literal style
}

\newtheorem{theo}{Théorème}[section]

\begin{document}

\begin{titlepage}

\newcommand{\HRule}{\rule{\linewidth}{0.5mm}} % Defines a new command for the horizontal lines, change thickness here

\begin{figure}
	\begin{minipage}[t]{.5\linewidth}
	\centering
	\includegraphics[scale=0.5]{esilv_logo.png}\\[1cm]
	\end{minipage}
	\begin{minipage}[t]{.5\linewidth}
	\includegraphics[scale=1.4]{edf_logo.png}\\[1cm]
	\end{minipage}
\end{figure}

\center % Center everything on the page


%	TITLE SECTION
%----------------------------------------------------------------------------------------

\HRule \\[0.4cm]
{ \huge \bfseries Pricer efficace d'actifs énergétiques}\\[0.4cm] % Title of your document
\HRule \\[1.5cm]

%----------------------------------------------------------------------------------------
%	AUTHOR SECTION
%----------------------------------------------------------------------------------------

\begin{minipage}{0.4\textwidth}
	\begin{flushleft} \large
		\emph{Team 16:}\\
		Anthony \textsc{Courillon} % Your name
		\\
		Denis \textsc{Barranco}
		\\
		Ousman \textsc{Asadullah}
		\\
		Hicham \textsc{Kortbi} % Your name
	\end{flushleft}
\end{minipage}
~
\begin{minipage}{0.4\textwidth}
	\begin{flushright} \large
		\emph{Partenaire:} \\
		Arnaud \textsc{de latour} % Supervisor's Name
		\\
		\emph{Mentor:} \\
		Martino \textsc{Grasselli} % Supervisor's Name
	\end{flushright}
\end{minipage}\\[2cm]

%----------------------------------------------------------------------------------------
%	DATE SECTION
%----------------------------------------------------------------------------------------

{\large Mars 2017}\\[2cm] % Date, change the \today to a set date if you want to be precise

\end{titlepage}


\chapter*{Introduction}
\addcontentsline{toc}{chapter}{Introduction}

Etant étudiants en quatrième année à l’ESILV dans la majeure finance et
dans le cadre du projet PI2, nous avons décidé de travailler sur la
réalisation d’un pricer d’actifs énergétiques proposé par notre partenaire EDF R\&D et encadré par Arnaud DE LATOUR

Un certain nombre d'actifs utilisés dans le domaine de l'énergie (par exemple les
centrales électriques) se prête facilement à une valorisation sous la forme de
produits dérivés complexes, de nature américaine (contrairement à une option
européenne, l'option américaine peut être exercée n'importe quand).

Le pricing de ces actifs met en oeuvre des méthodes numériques avancées qui
peuvent être coûteuses en temps de calcul. Lorsque ces méthodes sont fondées sur
une simulation Monte-Carlo (par exemple l'algorithme de Longstaff-Schwartz), on a
intérêt à employer des techniques de réduction de variance pour accélérer la
convergence de l'algorithme - et donc réduire le temps de calcul.

Cependant, l'efficacité de ces techniques dépend généralement du payoff de l'actif
qu'on cherche à valoriser. L'objectif de ce projet est donc d'implémenter et de
tester différentes méthodes de réduction de variance appliquées à des actifs
énergétiques, afin d'identifier les méthodes les plus pertinentes.

Nous allons réaliser ce pricer soous Python et nous allons utiliser exclusivement
ces 3 logiciels afin de réaliser ce projet:
\begin{itemize}
	\item Anaconda ;
	\item LiClipse ;
	\item Github.
\end{itemize}

\chapter{Evaluation du prix d'un produit dérivé par simulation}

\section{Principes des méthodes de Monte Carlo}

Le prix d'un produit dérivé d'un sous-jacent $X=(X_t)_{t\geq0}$ de payoff actualisé f(X), d'un marché financier sous probabilité risque neutre, s’écrit
comme l’espérance des flux futurs actualisés qu’il génère, $e_p$ :
\begin{equation}
	e_p = \mathbb{E}[f(X)]
\end{equation}
Il faut donc évaluer cette espérance pour connaître le prix d’un produit dérivé.
Nous allons utilize les méthodes de Monte Carlo qui consiste à simuler un grand
nombre de realisations, N indépendantes et identiquements distribuées, de f(X),
$f(X^i)$ avec $i\in \mathbb{N}$, puis on calcule chaque $f(X^i)$.

Enfin, on prend la moyenne empirique de ces réalisations, $\bar{e_N}$ estimateur de $e_p$ :
\begin{equation}
	\bar{e_N}=\frac{1}{N}\sum_{i=1}^{N} f(X^i) \xrightarrow[N\rightarrow{\infty}]{\text{p.s}} e_p \text{~avec~} \mathbb{E}[|f(X)|]<\infty
\end{equation}
la loi des grand nombres assure la convergence de $\bar{e_N}$ vers $e_p$.
Autrement dit, pour N assez grand, $\bar{e_N} \simeq e_p$.

\section{Limite pratique des méthodes de Monte Carlo}

\begin{theo}[Théorème Central Limite]\label{TCL}
	Nous allons utilisé le cas où on ne connait pas la loi de $e_p$.
	Si la loi de $e_p$ est inconnue (ou non normale), alors $\bar{e_N}\hookrightarrow N(\mu,\frac{\sigma}{\sqrt{N}})$ pour N assez grand, avec $\sigma$ = $\sqrt{(Var(f(X)))}$  et $\mu=e_p =\mathbb{E}[f(X)]$  alors on a:
	\begin{equation}
		\frac{\sqrt{N}(\bar{X}_n-\mu)}{\sigma}\hookrightarrow N(0,1)
	\end{equation}
\end{theo}

Donc le théorème central limite permet d'avoir l'approximation de l'erreur réalisée:
\begin{equation}
	\label{eq:approximation-erreur}
	\sqrt{N}(\frac{1}{N}\sum_{i=1}^{N}f(X^i)-e_p)\hookrightarrow N(0,\sqrt{Var(f(X))})
\end{equation}
ou :
\begin{equation}
	\sqrt{N}(\bar{e_N}-e_p)\hookrightarrow N(0,\sqrt{Var(f(X))})
\end{equation}

Cela montre le problème du nombre de simulations N nécessaires pour avoir un bon estimateur de $e_p$, plus la variance de $f(X^i)$ est élevé, plus N doit être grand.

Or N grand a pour conséquence une augmentation du temps de calcul. Donc plus $Var(f(X^i))$, plus l’obtention de l’évaluation de $e_p$ demandera un temps de calcul important.

Il est alors intéressant d’employer des techniques de réduction de variance pour accélérer la convergence de l’algorithme et donc réduire le temps de calcul. C’est l’un des problèmes que nous allons tenter de résoudre dans la réalisation de notre pricer.

Nous allons voir les différentes techniques de réduction de variance et tenter de les implémenter dans notre pricer dans le chapitre 2 de ce document.

\section{Contrôle de la convergence}

Il faut évaluer l’erreur d’approximation de \eqref{eq:approximation-erreur}  grâce à la construction d’un intervalle de confiance :
\begin{equation}
	Z = \frac{\sqrt{N}}{\sqrt{Var(f(X))}} (\frac{1}{N}\sum_{i=1}^{N}f(X^i)-e_p)\hookrightarrow N(0,1)
\end{equation}
En supposant que la variance est connue, on a :
\begin{equation}
	\mathbb{P}(-u\leq Z \leq u) = \alpha
\end{equation}
avec u fractile d'ordre 1-$\frac{\alpha}{2}$ de la N(0,1). Ce qui revient à :
\begin{equation}
	\mathbb{P}(-u \frac{\sqrt{Var(f(X))}}{\sqrt{N}}+\bar{e_N}\leq e_p \leq u \frac{\sqrt{Var(f(X))}}{\sqrt{N}}+\bar{e_N}) = \alpha
\end{equation}
et donc :
\begin{equation}
	\text{IC}_\alpha(e_p) = [-u \frac{\sqrt{Var(f(X))}}{\sqrt{N}}+\bar{e_N};u \frac{\sqrt{Var(f(X))}}{\sqrt{N}}+\bar{e_N}]
\end{equation}
L’erreur donné par un intervalle de confiance est :
\begin{equation}
	\mathbb{P}(e_p\in IC_\alpha ) \simeq \alpha
\end{equation}
Il a un problème c’est qu’on ne peut appliquer cette méthode que si la variance de f(X), Var(f(X)) est connue. Or ici, Var(f(X)) est inconnue, et il faut l’estimer grâce à un estimateur.

On peut construire un estimateur non biaisé tel que celui-ci :
\begin{equation}
	\bar{v_N} = \sqrt{\frac{1}{N-1}\sum_{i=1}^{N}(f(X^i)-\bar{e_N})^2}
\end{equation}

Voici l’implémentation sous Python de tout ce que nous avons énoncé dans cette partie I, on a pris  $\alpha$ =0.95 pour la contruction des intervalles de confiance.
\begin{verbatim}
MCResults = namedtuple('MCResults', ('mean', 'iclow', 'icup'))

def runmc(*allassets):
"""
Calcule par simulation Monte-Carlo les cash-flows des actifs de
`allassets`, sous la probabilité risque-neutre.
"""
allmarkets = set([m for asset in allassets for m in asset.getmarkets()])
marketsdates = set([d for market in allmarkets for d in market.getdates()])
assetsdates = set([d for asset in allassets for d in asset.getdates()])
earnings = {asset.name: {} for asset in allassets}
prices = {market.name: {} for market in allmarkets}
volumes = {asset.name: {m.name: {} for m in asset.getmarkets()} \
for asset in allassets}
alldates = assetsdates.union(marketsdates)
for date in alldates:
# Mise à jour des marchés :
for market in allmarkets:
market.simulate(date)
# Mise à jour des objets :
for asset in allassets:
earnings[asset.name][date] = asset.get_discounted_cf(date)
for market in asset.getmarkets():
volumes[asset.name][market.name][date] = asset.getvolume(date,
market)
if date not in prices[market.name]:
prices[market.name][date] = market.getspot(date)
earnings = pd.Panel(earnings).transpose(1, 0, 2)
volumes = {key: pd.Panel(value).transpose(1, 0, 2)
for key, value in volumes.items()}
prices = pd.Panel(prices).transpose(1, 0, 2)
return earnings, volumes, prices

def getmtm(cf, alpha=0.95):
"""
Calcule la MtM (i.e. : les cash-flows moyens réalisés et un intervalle
de confiance) de chaque actif listé dans `cf`.

Paramètres
----------
cf : pandas.Panel
Cash-flows réalisés pour chaque simulation (`items`), chaque
actif (`major_axis`) et chaque date (`minor_axis`).
alpha : double compris entre 0. et 1.
Quantile à utiliser pour le calcul des intervalles de confiances.
"""
nsims = cf.shape[0]
cumvalues = cf.sum(axis=2)
mean = cumvalues.mean(axis=1)
std = cumvalues.std(axis=1)
res = {}
for key, val in mean.items():
res[key] = mkmcresults(val, std[key], nsims)
return res

\end{verbatim}

\end{document}
